\clearpage
\section{The trace-zero Artin--Schreier tower}\label{app:odd_tower}


\hypertarget{odd_tower-exact-statement}{%
\subsection{Exact statement}\label{odd_tower-exact-statement}}

Let p be odd, F=F\_p((t)), and
0\textless h\_1\textless\ldots\textless h\_d with p not dividing any
h\_i. Put q\_i=(p+1)h\_i. The constructor below gives an actual matrix
B\_d over F of size p\^{}d with exactly the field, simplicity,
approximation-level, odd-prime quotient, and conductor conclusions of
\path{RECURSIVE_TOWER166.md}. In addition, its matrix trace is exactly zero.
Every native shorter constructor and every amplified arithmetic
approximation also has matrix trace zero. Therefore every difference
between consecutive arithmetic approximations has matrix trace zero.

The trace statement is an identity over the current coefficient field,
not merely a congruence, a field trace inferred from characteristic
zero, or a division by the matrix size. All coefficient fields have
characteristic p, so division by p\^{}j would be invalid.

\hypertarget{odd_tower-an-explicit-separable-minimal-extension-with-a-traceless-parameter}{%
\subsection{1. An explicit separable minimal extension with a traceless
parameter}\label{odd_tower-an-explicit-separable-minimal-extension-with-a-traceless-parameter}}

At an arbitrary recursive step let L have residue field F\_p and
uniformizer t\_L. Given q\textgreater0 prime to p, set

\begin{Verbatim}[fontsize=\small,breaklines=true,breakanywhere=true]
E=L(alpha),  alpha^p-alpha=t_L^(-q).
\end{Verbatim}

The right side is not an Artin--Schreier coboundary in L: the negative
valuation of x\^{}p-x, if x has negative valuation, is divisible by p.
Thus the polynomial is irreducible of degree p.~The extension is
separable, and the valuation equation forces ramification index p and

\begin{Verbatim}[fontsize=\small,breaklines=true,breakanywhere=true]
v_E(alpha)=-q.
\end{Verbatim}

It is therefore totally ramified and alpha is minimal over L. For odd p
its field trace is

\begin{Verbatim}[fontsize=\small,breaklines=true,breakanywhere=true]
Tr_(E/L)(alpha)=sum_(c in F_p)(alpha+c)=0.
\end{Verbatim}

Choose the unique b in \{1,\ldots,p-1\} with qb congruent to -1 modulo
p, put a=(qb+1)/p, and define

\begin{Verbatim}[fontsize=\small,breaklines=true,breakanywhere=true]
z=t_L^a alpha^b.
\end{Verbatim}

Then v\_E(z)=pa-qb=1, so z is a uniformizer and E=L(z). Its monic
minimal polynomial \(f(X)=X^p+a_{p-1}X^{p-1}+\cdots+a_0\)
is Eisenstein. It is an explicit polynomial: compute the product over
the p conjugates

\begin{Verbatim}[fontsize=\small,breaklines=true,breakanywhere=true]
f(X)=product_(c in F_p) (X-t_L^a(alpha+c)^b)
\end{Verbatim}

and reduce coefficients using \(\alpha^p-\alpha=t_L^{-q}\).
The invariant coefficients belong to L. This is a finite algebraic
instruction, with no choice of an unknown extension or matrix.

\hypertarget{odd_tower-normalized-corestriction-and-an-explicit-filtered-right-inverse}{%
\subsection{2. Normalized corestriction and an explicit filtered right
inverse}\label{odd_tower-normalized-corestriction-and-an-explicit-filtered-right-inverse}}

Use the basis 1,z,\ldots,z\^{}(p-1) of E over L. If T\_ij is its
coefficient matrix unit and pr is the trace corestriction, then

\begin{Verbatim}[fontsize=\small,breaklines=true,breakanywhere=true]
pr(T_ij)=z^i z_j^*,
z_j^*=sum_(k=j+1)^p a_k z^(k-j-1) / f'(z),  a_p=1.
\end{Verbatim}

Set

\begin{Verbatim}[fontsize=\small,breaklines=true,breakanywhere=true]
kappa=-f'(z)/z^(p-1),  s=kappa pr.
\end{Verbatim}

The displayed dual-basis formula gives

\begin{Verbatim}[fontsize=\small,breaklines=true,breakanywhere=true]
s(T_ij)=-z^(i-j) + terms of strictly larger E-valuation.
\end{Verbatim}

Indeed v\_E(a\_k)\textgreater=p for 1\textless=k\textless p, so the
contribution indexed by a\_k has valuation at least k greater than the
first contribution. Applying s entrywise on E\^{}m gives a B-bimodule
corestriction from End\_L(E\^{}m) to B=End\_E(E\^{}m). With the complete
interleaved period-pm L-chain and period-m E-chain, it preserves grades
and satisfies \(s(P^r)=Q^r\).

Let P0 extract the constant coefficient on every E-coordinate. The
constant dual coefficient satisfies

\begin{Verbatim}[fontsize=\small,breaklines=true,breakanywhere=true]
z_0^*=-a_0/(z f'(z)),
u_0=s(P0)=kappa z_0^*=a_0/z^p in O_E^times.
\end{Verbatim}

For any C in M\_m(E), define the actual L-linear matrix

\begin{Verbatim}[fontsize=\small,breaklines=true,breakanywhere=true]
Lift_E(C)=u_0^(-1) C P0.
\end{Verbatim}

Then

\begin{Verbatim}[fontsize=\small,breaklines=true,breakanywhere=true]
s(Lift_E(C))=C.
\end{Verbatim}

This lift preserves the exact filtration valuation. To see this, first
multiply C by the unit u\_0\^{}(-1), which does not change its valuation
in the E-chain. Expand each entry in the basis 1,z,\ldots,z\^{}(p-1).
The only nonzero columns of Lift are its constant-coefficient columns.
The L-grade of a coefficient c\_i z\^{}i E\_uv is

\begin{Verbatim}[fontsize=\small,breaklines=true,breakanywhere=true]
pm*v_L(c_i)+mi+u-v,
\end{Verbatim}

equal to its E-grade, and distinct i have valuations distinct modulo p.
Thus the minimum cannot cancel. This also proves the stated surjectivity
\(s(P^r)=Q^r\) without a compactness or lifting-existence
assumption.

\hypertarget{odd_tower-exact-trace-identity}{%
\subsection{3. Exact trace identity}\label{odd_tower-exact-trace-identity}}

For y in E, extraction of its constant coefficient is

\begin{Verbatim}[fontsize=\small,breaklines=true,breakanywhere=true]
coeff_0(y)=Tr_(E/L)(z_0^* y).
\end{Verbatim}

The ordinary L-matrix trace of the rank-one operator yP0 is coeff\_0(y).
Consequently, for every C in M\_m(E),

\begin{Verbatim}[fontsize=\small,breaklines=true,breakanywhere=true]
tr_L(Lift_E(C))
   =Tr_(E/L)(z_0^* u_0^(-1) tr_E(C))
   =Tr_(E/L)(kappa^(-1) tr_E(C)).             (T)
\end{Verbatim}

Here tr\_E(C) is the ordinary m-by-m matrix trace over E. In particular,
if tr\_E(C)=0, then the lift has ordinary L-matrix trace zero exactly.

\hypertarget{odd_tower-the-recursive-actual-matrix}{%
\subsection{4. The recursive actual
matrix}\label{odd_tower-the-recursive-actual-matrix}}

Define B(L,t\_L;h\_1,\ldots,h\_d) by induction. For d=1 use the
multiplication matrix of the alpha from Section~1 with q=q\_1. For
d\textgreater=2 use Section~1 with q=q\_d, put m=p\^{}(d-1), recursively
construct

\begin{Verbatim}[fontsize=\small,breaklines=true,breakanywhere=true]
C=B(E,z;h_1,...,h_(d-1)),
\end{Verbatim}

and set

\begin{Verbatim}[fontsize=\small,breaklines=true,breakanywhere=true]
B(L,t_L;h_1,...,h_d)=alpha I_m+Lift_E(C).       (ASR)
\end{Verbatim}

Everything in (ASR) is explicit in finite field operations and the
displayed degree-p polynomial extensions. It differs from the old
conjugacy-preserving two-summand construction. Preserving old lower
summand conjugacy classes is not claimed.

Inductively tr\_E(C)=0. Equation (T) and the Artin--Schreier trace give

\begin{Verbatim}[fontsize=\small,breaklines=true,breakanywhere=true]
tr_L(B)=m Tr_(E/L)(alpha)+tr_L(Lift_E(C))=0.
\end{Verbatim}

For the induction base the trace is Tr\_(E/L)(alpha)=0. Thus this
argument proves tracelessness in the current base field for every native
rank, not merely tracelessness after an amplification whose size is
divisible by p.

\hypertarget{odd_tower-actual-fields-and-profiles}{%
\subsection{5. Actual fields and
profiles}\label{odd_tower-actual-fields-and-profiles}}

Set gamma=alpha I\_m and b=Lift\_E(C). The minimal outer stratum has

\begin{Verbatim}[fontsize=\small,breaklines=true,breakanywhere=true]
n=p^(d-1)q_d,  k_0(gamma)=-n.
\end{Verbatim}

The inner actual maximal field supplied by induction has

\begin{Verbatim}[fontsize=\small,breaklines=true,breakanywhere=true]
R=p^(d-2)q_(d-1),  v(C)=-R,  k_0(C)=-q_1.
\end{Verbatim}

We have s(b)=C, v(b)=-R, and n\textgreater pR\textgreater2R because p is
odd and the h\_i increase. Hence

\begin{Verbatim}[fontsize=\small,breaklines=true,breakanywhere=true]
-2R > -n-q_1 = k_0(gamma)+k_0(C).
\end{Verbatim}

BK Remark 2.2.7 applies to this pure derived maximal field. It gives the
field generated by the actual displayed gamma+b, with degree and
ramification p\^{}d, residue degree one, normalization of the complete
chain, and k\_0=-q\_1. Its valuation remains -n.~Thus {[}A,n,0,B{]} is
simple.

This repeats the precise argument of \path{RECURSIVE_TOWER166.md} with a new
outer minimal extension and normalized right inverse. No step of that
argument required the old polynomial e\^{}p+t e+t.

\hypertarget{odd_tower-native-approximations-amplification-and-traceless-corrections}{%
\subsection{6. Native approximations, amplification, and traceless
corrections}\label{odd_tower-native-approximations-amplification-and-traceless-corrections}}

For the full rank p\^{}d matrix put, for 1\textless=j\textless=d,

\begin{Verbatim}[fontsize=\small,breaklines=true,breakanywhere=true]
B^[j]=B(F,t;h_(d-j+1),...,h_d) tensor I_(p^(d-j)).
\end{Verbatim}

These are literal truncations of (ASR). At the outer level the same
largest q\_d selects the same Artin--Schreier extension and uniformizer
in all truncations. At the next level the same q\_(d-1) does the same,
and so on. Thus no hidden change of bases or incompatible parameter
fields enters the truncations. Lift\_E commutes with amplification by an
identity because P0 acts separately on each coordinate.

Every native shorter B has trace zero by Section~4. Its amplification
therefore has trace zero. In particular

\begin{Verbatim}[fontsize=\small,breaklines=true,breakanywhere=true]
tr_F(B^[j])=0,
tr_F(B^[j+1]-B^[j])=0.                        (Z)
\end{Verbatim}

The same is true at native rank p\^{}(j+1): the correction between
B(F,t;h\_(d-j),\ldots,h\_d) and the identity-p amplification of
B(F,t;h\_(d-j+1),\ldots,h\_d) has trace zero. This native assertion is
the one relevant when constructing characters before each transfer.

The field degrees and critical levels of \(B^{[j]}\) are exactly

\begin{Verbatim}[fontsize=\small,breaklines=true,breakanywhere=true]
degree p^j,  n_j=p^(d-j)q_(d-j+1).
\end{Verbatim}

Valuation-preservation of every Lift gives

\begin{Verbatim}[fontsize=\small,breaklines=true,breakanywhere=true]
v_A(B^[j+1]-B^[j])=-n_(j+1).
\end{Verbatim}

Thus the approximation and defining-sequence proof of the first note
applies unchanged. The odd-prime J/H dimensions, grade-index formula,
Swan conductor and self-pair conductor depend on these same levels and
centralizer degrees, so all their formulas carry over unchanged.

\hypertarget{odd_tower-precisely-what-the-trace-supplies}{%
\subsection{7. Precisely what the trace
supplies}\label{odd_tower-precisely-what-the-trace-supplies}}

For any displayed correction delta and any scalar u in L\^{}times for
which its additive correction character is defined,

\begin{Verbatim}[fontsize=\small,breaklines=true,breakanywhere=true]
psi_delta(u I)=psi_L((u-1)tr_L(delta))=1.
\end{Verbatim}

This proves triviality of the actual additive correction on the scalar
subgroup of its domain. It does not, by itself, construct a simple
character with prescribed values on larger scalar groups. Compatible
extension, transfer, and intersection arguments for the full simple
character are separate obligations.

The note also does not identify the constructed candidates with any Weil
source, choose them by inverse LLC, realize arbitrary stabilizer
operators, or prove same-candidate H5. Its new output is an explicit
arithmetic defining sequence whose trace vanishes at every native and
amplified stage, suitable as input to that separate character argument.

\hypertarget{odd_tower-primary-dependency}{%
\subsection{Primary dependency}\label{odd_tower-primary-dependency}}

Bushnell--Kutzko, The admissible dual of GL(N) via compact open
subgroups, Remark 2.2.7 (printed p. 55), with the defining-sequence and
ring recursion already pinned in \path{RECURSIVE_TOWER166.md}. Retained source
SHA256:
\path{d69db0827256707a34cd2562af68b9de11fc889ab3f43535ec47b620d58ca9c8}. The
trace-dual and Artin--Schreier calculations above are given in full.
